\documentclass[12pt,a4paper]{article}

\usepackage[utf8]{inputenc}
\usepackage[T1]{fontenc}
\usepackage[french]{babel}

\usepackage{styles/couleurs}
\usepackage{styles/mise_en_page}
\usepackage{styles/boites}
\usepackage{styles/maths}
\usepackage{styles/figures}

\usepackage{amsmath}
\usepackage{amssymb}
\usepackage{enumitem}
\usepackage{array}
\usepackage{tabularx}
\usepackage{graphicx}
\usepackage{tikz}
\usepackage{xparse}
\usetikzlibrary{calc,angles,quotes}

\setlist[itemize]{leftmargin=1.4em,itemsep=0.15em,topsep=0.2em}
\setlist[enumerate]{leftmargin=*,itemsep=0.22em,topsep=0.22em}

\newcommand{\titresequence}[1]{%
\begin{center}
  {\sffamily\LARGE\bfseries\color{TitreC}Séquence 2 -- Théorème de Pythagore}\\[0.2em]
  {\sffamily\large\bfseries\color{SavoirC}#1}
\end{center}}

\newcommand{\lignepointillee}{\noindent\dotfill\par\vspace{0.45em}}
\newcommand{\lignesreponse}[1]{%
\begin{tcolorbox}[
  enhanced,colback=white,colframe=NeutreC!40,boxrule=0.5pt,arc=1.5mm,
  left=3mm,right=3mm,top=2mm,bottom=1mm,before skip=2pt,after skip=2pt]
  \foreach \i in {1,...,#1}{\lignepointillee}
\end{tcolorbox}}

\begin{document}
\titresequence{Activité 4 -- Deux calculs, une rédaction}

\section*{Cas 1 -- Calculer l'hypoténuse}

Le triangle $ABC$ est rectangle en $A$, avec $AB=4$ cm et $AC=7$ cm.

\begin{enumerate}
  \item Quel côté cherche-t-on ? Est-ce l'hypoténuse ?
  \item Écris l'égalité de Pythagore.
  \item Calcule $BC^2$.
  \item Donne la valeur exacte de $BC$, puis une valeur approchée au dixième.
\end{enumerate}
\lignesreponse{8}

\section*{Cas 2 -- Calculer un côté de l'angle droit}

Le triangle $DEF$ est rectangle en $E$, avec $DF=8$ cm et $EF=6$ cm.

\begin{enumerate}
  \item Quel côté est l'hypoténuse ?
  \item Écris l'égalité de Pythagore sans transformer immédiatement la formule.
  \item Isole ensuite $DE^2$.
  \item Donne la valeur exacte de $DE$, puis une valeur approchée au dixième.
\end{enumerate}
\lignesreponse{9}

\section*{Comparer les deux méthodes}

Pourquoi additionne-t-on dans le premier calcul et soustrait-on dans le second ?
\lignesreponse{4}

\begin{tcolorbox}[colback=SavoirC!5,colframe=SavoirC!60,boxrule=0.6pt,arc=2mm]
\textbf{Trame de rédaction.}
\begin{enumerate}
  \item On sait que le triangle est rectangle en \ldots
  \item D'après le théorème de Pythagore : \ldots
  \item On remplace par les longueurs connues et on calcule.
  \item On conclut avec une unité et, si nécessaire, un arrondi.
\end{enumerate}
\end{tcolorbox}
\end{document}
